% !TeX program = XeLaTeX
% !TeX root = ../../shloka.tex
% --meta--
% deity: Ganga
% composer: Adi Shankaracharya
% type: Stotram
% wiki_title: Ganga Stotram (Shankara)
% --end-meta--
% Text: श्रीगङ्गाकल्पोक्तपूजाविधिः (Idkidu, 2015), pp. 51-52.

\sect{गङ्गास्तोत्रम्}

\twolineshloka
{देवि सुरेश्वरि भगवति गङ्गे त्रिभुवनतारिणि तरलतरङ्गे}
{शङ्करमौलिविहारिणि विमले मम मतिरास्तां तव पदकमले} %॥१॥

\twolineshloka
{भागीरथि सुखदायिनि मातस्तव जलमहिमा निगमे ख्यातः}
{नाहं जाने तव महिमानं पाहि कृपामयि मामज्ञानम्} %॥२॥

\twolineshloka
{हरिपदपाद्यतरङ्गिणि गङ्गे हिमविधुमुक्ताधवलतरङ्गे}
{दूरीकुरु मम दुष्कृतिभारं कुरु कृपया भवसागरपारम्} %॥३॥

\twolineshloka
{तव जलममलं येन निपीतं परमपदं खलु तेन गृहीतम्}
{मातर्गङ्गे त्वयि यो भक्तः किल तं द्रष्टुं न यमः शक्तः} %॥४॥

\twolineshloka
{पतितोद्धारिणि जाह्नवि गङ्गे खण्डितगिरिवरमण्डितभङ्गे}
{भीष्मजननि हे मुनिवरकन्ये पतितनिवारिणि त्रिभुवनधन्ये} %॥५॥

\twolineshloka
{कल्पलतामिव फलदां लोके प्रणमति यस्त्वां न पतति शोके}
{पारावारविहारिणि गङ्गे विमुखयुवतिकृततरलापाङ्गे} %॥६॥

\twolineshloka
{तव चेन्मातः स्रोतःस्नातः पुनरपि जठरे सोऽपि न जातः}
{नरकनिवारिणि जाह्नवि गङ्गे कलुषविनाशिनि महिमोत्तुङ्गे} %॥७॥

\twolineshloka
{पुनरसदङ्गे पुण्यतरङ्गे जय जय जाह्नवि करुणापाङ्गे}
{इन्द्रमुकुटमणिराजितचरणे सुखदे शुभदे भृत्यशरण्ये} %॥८॥

\twolineshloka
{रोगं शोकं पापं तापं हर मे भगवति कुमतिकलापम्}
{त्रिभुवनसारे वसुधाहारे त्वमसि गतिर्मम खलु संसारे} %॥९॥

\twolineshloka
{अलकानन्दे परमानन्दे कुरु मयि करुणां कातरवन्द्ये}
{तव तटनिकटे यस्य हि वासः खलु वैकुण्ठे तस्य निवासः} %॥१०॥

\twolineshloka
{वरमिह नीरे कमठो मीनः किं वा तीरे शरटः क्षीणः}
{अथवा श्वपचो मलिनो दीनस्तव न हि दूरे नृपतिकुलीनः} %॥११॥

\twolineshloka
{भो भुवनेश्वरि पुण्ये धन्ये देवि द्रवमयि मुनिवरकन्ये}
{गङ्गास्तवमिमममलं नित्यं पठति नरो यः स जयति सत्यम्} %॥१२॥

\twolineshloka
{येषां हृदये गङ्गाभक्तिस्तेषां भवति सदा सुखमुक्तिः}
{मधुराकान्तापज्झटिकाभिः परमानन्दकलितललिताभिः} %॥१३॥

\twolineshloka
{गङ्गास्तोत्रमिदं भवसारं वाञ्छितफलदं विमलं सारम्}
{शङ्करसेवकशङ्कररचितं पठति सुखी स्तव इति च समाप्तः} %॥१४॥

{॥इति श्रीमच्छङ्कराचार्यविरचितं गङ्गास्तोत्रं सम्पूर्णम्॥}
